\documentclass[10pt,a4paper]{amsart}
\usepackage[margin=19mm]{geometry}
\usepackage{amsmath,amssymb,mathtools}
\usepackage[scr=rsfs,cal=euler]{mathalfa}
\usepackage{xfrac}
\usepackage[unicode,hidelinks,bookmarks=false]{hyperref}
\newcommand{\ie}{i.e.\ }
\newcommand{\cf}{cf.\ }
\newcommand{\resp}{respectively\ }
\newcommand{\Subsection}{\subsection}
\providecommand{\nobreakdash}{\nobreak\hbox{-}}
\setlength{\emergencystretch}{2em}
\multlinegap0pt
\usepackage{bm}
\usepackage{bbm}
\usepackage{dsfont}
\usepackage{xspace}
\usepackage{cancel}
\usepackage{mathtools}
\usepackage{amsthm}
\makeatletter
\@ifundefined{lemma}{\newtheorem{lemma}{Lemma}}{}
\makeatother
\title[Morrey's Derivation of Hydrodynamics from Statistical Mechan]{Morrey's Derivation of Hydrodynamics from Statistical Mechanics: A Modern Exposition}
\author{Yuyang Wang}
\date{Source: arXiv:2608.10612v1 (CC BY 4.0)}
\begin{document}
\maketitle

\noindent\emph{Source and licence.} Morrey's Derivation of Hydrodynamics from Statistical Mechanics: A Modern Exposition. Version arXiv:2608.10612v1 (CC BY 4.0), distributed under the Creative Commons Attribution 4.0 International licence, as recorded in the arXiv licence metadata for this exact version and shown on the abstract page https://arxiv.org/abs/2608.10612v1. Full rights notice: licenses/arxiv-2608.10612-CC-BY-4.0.txt.\par\smallskip

\noindent\textit{Editorial scope.} Counted unit: the Lemma whose source label is \texttt{lem:bounds} in the article (source0.tex lines 1761--1807, source label \texttt{lem:bounds}), delivered together with the article's own complete original proof of it (source0.tex lines 1808--1985, one \texttt{proof} environment of 178 lines). No mathematics is written, repaired or re-derived: statement and proof are the article's own text line for line. Required context retained uncounted: none. Every definition, display and label of those lines is retained unchanged. Omitted from this package and not redistributed: 1--1760, 1986--3450. Nothing inside a retained excerpt is omitted or altered: every retained body is delivered exactly as the article's lines are written, so no omission receipt of any kind accompanies this package. Citations: the retained excerpts carry no citation command at all, so no bibliography entry of the article is transcribed and no reference is redistributed. Reference adaptation: none was needed and none was made; every reference inside the delivered unit resolves inside it, so no unresolved reference or citation is produced. Printed slips kept literal: no printed word, symbol, hypothesis or number of the counted statement or of its proof was altered. Layout and class: the article's own class is \texttt{article} with packages including amsthm, amsmath, amssymb, amsfonts, physics, geometry, hyperref, fancyhdr, titlesec, enumitem, graphicx, xcolor, caption, tikz; the wrapper is the lane's pinned \texttt{amsart} editorial wrapper at 10pt on A4, which restates the theorem-like environments the retained text needs and adds the article's own macro definitions that the retained text needs (none), with extra wrapper packages (bm, bbm, dsfont, xspace, cancel, mathtools). The delivered numbering is the unit's own. Assets: no retained span contains a figure environment, a graphics-inclusion command, a PDF-page inclusion, a table or a TikZ picture; no image file is copied into this package, the package declares no asset dependency and no image file is redistributed with it. Licence: version arXiv:2608.10612v1 of this article is distributed under the Creative Commons Attribution 4.0 International licence, as recorded in the arXiv licence metadata for this exact version and shown on the abstract page https://arxiv.org/abs/2608.10612v1. A full rights notice is delivered at licenses/arxiv-2608.10612-CC-BY-4.0.txt. Characters: every retained excerpt is pure ASCII, so the delivered engine, XeLaTeX with its default Latin Modern text font, reports no missing character.
\begin{lemma}[6.2]\label{lem:bounds}
    For all integers $n,k,l$ and points $x_1,\ldots,x_l \in R$, let $m = -\min \Phi(r).$
Then, for $\rho$ small enough,
\begin{equation}
\frac{1}{12}\left[1-\frac{2}{3}\phi(h)\right]^n
\le
D^{\,n-k}_{N kl}(x_1,\ldots,x_l;\rho,b)
\le
\left(1+\frac{2}{3}bKL\right)^{n-k}
\exp\!\left\{\frac{l(l-1)mKb}{1+2bKL/3}\right\},
\tag{6.6}\label{6.6}
\end{equation}
and
\begin{equation}
\frac{1}{12}\left[1-\frac{2}{3}\phi(h)\right]^n
\le
C^{\,n-k}_{N kl}(\rho,b)
\le
\left(1+\frac{2}{3}bKL\right)^{n-k},
\qquad N>l .
\tag{6.7}\label{6.7}
\end{equation}
Here
\[
h=\frac{\rho \tau_0(b)}{D},
\qquad
\phi(h)=\sum_{\nu=1}^{\infty}\frac{h^\nu}{\nu(\nu+1)},
\qquad
0\le h\le1 .
\]
If $\Phi(r)\le0$ for $r\ge r_0$, then
\begin{equation}
C^{\,n-k}_{N kl}(\rho,b)\ge e^{-n\gamma},
\qquad
D^{\,n-k}_{N kl}\ge e^{-n\gamma},
\tag{6.8}\label{6.8}
\end{equation}
where
\[
\gamma :=\tfrac23\phi(h_0),
\qquad
h_0:=\frac{4\pi r_0^3\rho}{3D},
\qquad
h_0\le1 .
\]

\end{lemma}

\begin{proof}
We first prove the lower bound in \eqref{6.6}. As in Morrey, write
\[
\Phi=\Phi_0-\Phi_1
\]
with the notation from Section 8. In particular,
\[
W^{**}_{0pl}=W^{***}_{0Nl}\big|_{N\mapsto p},\qquad W^{**}_{0ll}=0.
\]
Then for $l\le p < N$ we have
\begin{equation}
W^{**}_{0,p+1,l}+W^{*}_{0,p+1,l}
=
2K\sum_{i=1}^p \Phi_0\!\left(\frac{|x_i-x_{p+1}|}{\epsilon}\right)
+
W^{**}_{0pl}+W^{*}_{0pl}.
\tag{6.9}\label{6.9}
\end{equation}

Now recall that in our notation
\[
ND\epsilon^3=pr.
\]
Hence
\begin{align}
&r^{-1}\int_R
\left[
1-\exp\left\{-2Kb\sum_{i=1}^p \Phi_0\!\left(\frac{|x_i-x_{p+1}|}{\epsilon}\right)\right\}
\right]dx_{p+1}
\notag\\
&\le
\sum_{i=1}^p
r^{-1}\int_R
\left[
1-\exp\left\{-2Kb\,\Phi_0\!\left(\frac{|x_i-x_{p+1}|}{\epsilon}\right)\right\}
\right]dx_{p+1}
\notag\\
&\le
\frac{p}{N}\cdot \frac{\rho}{D}
\int_{\mathbb R^3}
\left[1-\exp\{-2Kb\,\Phi_0(|\xi|)\}\right]\,d\xi
=
\frac{p}{N}h.
\tag{6.10}\label{6.10}
\end{align}
Here $h$ is the quantity introduced earlier in the statement of the lemma.

Using \eqref{6.9} and \eqref{6.10} with $p=N-1$, we obtain
\begin{align*}
&r^{-N+l}\int_{R\times\cdots\times R}
\exp\{-bW^{**}_{0Nl}-bW^{*}_{0Nl}\}\,dx_{l+1}\cdots dx_N
\\
&=
r^{-N+l+1}\int_{R\times\cdots\times R}
\Bigg[
r^{-1}\int_R
\exp\left\{-2Kb\sum_{i=1}^{N-1}\Phi_0\!\left(\frac{|x_i-x_N|}{\epsilon}\right)\right\}dx_N
\Bigg]
\\
&\hspace{4cm}\times
\exp\{-bW^{**}_{0,N-1,l}-bW^{*}_{0,N-1,l}\}\,dx_{l+1}\cdots dx_{N-1}
\\
&\ge
\left[1-\frac{N-1}{N}h\right]
r^{-N+l+1}
\int_{R\times\cdots\times R}
\exp\{-bW^{**}_{0,N-1,l}-bW^{*}_{0,N-1,l}\}\,dx_{l+1}\cdots dx_{N-1}.
\end{align*}
Repeating this argument successively for $x_N,x_{N-1},\dots,x_{l+1}$ gives
\begin{equation}
r^{-N+l}\int_{R\times\cdots\times R}
\exp\{-bW^{**}_{0Nl}-bW^{*}_{0Nl}\}\,dx_{l+1}\cdots dx_N
\ge
\prod_{j=l}^{N-1}\left(1-\frac{j}{N}h\right)
=: \Pi_N(h).
\tag{6.11}\label{6.11}
\end{equation}

Next, estimate $\Pi_N(h)$ from below. We write
\begin{equation*}
-\log \Pi_N(h)
=
\sum_{\nu=1}^\infty \frac{h^\nu}{\nu}\sum_{j=0}^{N-1}\left(\frac{j}{N}\right)^\nu
\le N\phi(h),
\tag{6.12}\label{6.12}
\end{equation*}
so that
\[
\Pi_N(h)\ge e^{-N\phi(h)}=:e^{-n\beta},
\qquad
\beta=\frac{2}{3}\phi(h).
\]

Now choose $s$ so that
\[
\beta<bs<1,
\]
and let $\Sigma$ denote the set of $(x_{l+1},\dots,x_N)$ for which
\[
n^{-1}(W^{**}_{0Nl}+W^{*}_{0Nl})\le s.
\]
Then
\begin{align*}
D^{\,n-k}_{Nkl}
&\ge D^{\,n-k}_{0Nkl}
\notag\\
&\ge
r^{-N+l}\int_\Sigma
\left[
\exp\{bW^{**}_{0Nl}+bW^{*}_{0Nl}\}
\cdot
\Bigl(\mathrm{pos}\bigl(1-\tfrac{b}{n}W^{**}_{0Nl}-\tfrac{b}{n}W^{*}_{0Nl}\bigr)\Bigr)^{n-k}
\right]
\notag\\
&\times
\exp\{-bW^{**}_{0Nl}-bW^{*}_{0Nl}\}\,dx_{l+1}\cdots dx_N\\
&\ge
\bigl[e^{bs}(1-bs)\bigr]^n
\, r^{-N+l}\int_\Sigma
\exp\{-bW^{**}_{0Nl}-bW^{*}_{0Nl}\}\,dx_{l+1}\cdots dx_N
\notag\\
&\ge
\bigl(e^{-nbs-n\beta}-1\bigr)(1-bs)^n.
\tag{6.13}\label{6.13}
\end{align*}

On the right-hand side, set
\[
bs=\beta+\frac{y}{n}.
\]
Then \eqref{6.13} becomes
\begin{equation}
(1-\beta)^n(e^y-1)\left[1-\frac{1}{n}\frac{y}{1-\beta}\right]^n,
\qquad
0<y<n(1-\beta).
\tag{6.14}\label{6.14}
\end{equation}
For fixed $y$, the third factor increases with $n$, so \eqref{6.14} is bounded below by
\[
(1-\beta)^n\left[y-\frac{y^2}{1-\beta}\right],
\qquad
0<y<1-\beta.
\]
This gives the desired left-hand inequality in \eqref{6.6}.

The left-hand inequality in \eqref{6.7} for $l=0$ is just a special case of \eqref{6.6}. For general $l$, we use
\begin{equation}
C_{Nkl}(\rho,b)
=
C_{N-l,\,k-3l/2,\,0}\bigl(\rho(1-l/N),b\bigr),
\tag{6.15}\label{6.15}
\end{equation}
since
\[
(N-l)D\epsilon^3=(1-l/N)r
\]
in our notation. Hence the same lower bound follows for all $l$.

For the remaining lower bound involving the $D$'s, let $\Sigma_0\subset R^{N-l}$ be the set on which
\[
\frac{|x_j-x_k|}{\epsilon}\ge r_0
\]
for all $j,k=l+1,\dots,N$, and also for $j=1,\dots,l$, $k=l+1,\dots,N$. By choosing the variables one by one, we obtain
\begin{equation}
r^{-N+l}\mu(\Sigma_0)
\ge
\left(1-\frac{lh_0}{N}\right)\cdots
\left(1-\frac{(N-1)h_0}{N}\right)
\ge
\Pi_N(h_0),
\tag{6.16}\label{6.16}
\end{equation}
where
\[
h_0=\frac{4}{3}\pi r_0^3\epsilon^3r^{-1}N=\frac{4\pi r_0^3\rho}{3D}.
\]
Now the conclusion follows from \eqref{6.12}, since on $\Sigma_0$ the integrand in $D^{\,n-k}_{Nk}$ is bounded below by $1$.
\end{proof}

\end{document}
